\chapter{Installation et lancement}

\section{Prérequis}

Pour une installation depuis les sources, les prérequis recommandés sont :

\begin{itemize}
  \item Windows 10 ou Windows 11 ;
  \item Python 3.12 ;
  \item un environnement virtuel Python ;
  \item les dépendances listées dans \fichier{requirements.txt}.
\end{itemize}

Le build Windows publié cible Windows 10/11. Les anciennes versions de Windows
ne sont pas supportées par le couple Python 3.12 et Qt 6 utilisé par
l'application.

\section{Installation depuis les sources}

\begin{verbatim}
git clone https://github.com/chakibhenaoui/HEXA_Structures.git
cd HEXA_Structures

py -3.12 -m venv .venv
.\.venv\Scripts\Activate.ps1
.venv\Scripts\activate.bat

pip install -r requirements.txt
python main.py
\end{verbatim}

\section{OpenSeesPy optionnel}

OpenSeesPy n'est pas nécessaire pour lancer le logiciel avec le moteur PyNite.
Pour utiliser le backend OpenSeesPy, il doit être installé séparément :

\begin{verbatim}
pip install openseespy
\end{verbatim}

Si OpenSeesPy est installé dans un dossier spécifique, la variable
d'environnement \path{HEXA_PYTHON_SITE_PACKAGES} peut pointer vers le dossier
\fichier{site-packages} concerné.

\section{Construction de l'exécutable Windows}

\begin{verbatim}
pip install pyinstaller
build.bat
\end{verbatim}

Le résultat attendu est :

\begin{verbatim}
dist\HEXA Structures\HEXA Structures.exe
\end{verbatim}

\section{Vérification rapide}

Après lancement, vérifier que :

\begin{itemize}
  \item la fenêtre principale s'ouvre sans erreur ;
  \item l'arbre du modèle est visible ;
  \item le panneau de propriétés est visible ;
  \item la vue graphique est disponible ou affiche un message clair si la vue
  3D ne peut pas être initialisée ;
  \item le menu \menu{Analyse} permet de choisir le moteur de calcul.
\end{itemize}

\captureplaceholder{Fenêtre principale au premier lancement}{Écran d'accueil du logiciel}
